% Full editorial derivations for ED26-BM-RAT2021-10-4.
\section{Rational-choice social explanation: provision-game derivations}
\label{annex:p17-rat104-provision-game}

\begin{annexitem}{annex:p17.rat104-vod-equilibria}{thm:p17-rat104-vod-equilibria}{Pure equilibria, symmetric mixing and the macro effect}
\begin{proof}
At a pure profile with no volunteer, every actor obtains zero and can
obtain \(U-K>0\) by volunteering. Such a profile is not an equilibrium.
If at least two actors volunteer, a volunteer can abstain while the good
remains supplied, raising that actor's utility from \(U-K\) to \(U\).
Such a profile is not an equilibrium either. If exactly one volunteers,
that actor would reduce utility from \(U-K\) to zero by abstaining;
each other actor would reduce utility from \(U\) to \(U-K\) by
volunteering. Hence exactly the \(N\) sole-volunteer profiles are pure
equilibria, and they all guarantee provision.

Now suppose the \(N-1\) opponents of an actor contribute independently
with common probability \(p\). The payoff from contributing is \(U-K\).
The payoff from abstaining is
\[
                   U\bigl[1-(1-p)^{N-1}\bigr].
\]
In a symmetric equilibrium \(p=0\) is impossible by the first pure
argument, and \(p=1\) is impossible by the second. An interior mixture
can be a best response only when the two pure actions have equal expected
utility, giving
\[
            U(1-p)^{N-1}=K,
\qquad p=1-(K/U)^{1/(N-1)}=p^*.
\]
Because \(0<K/U<1\), this equation has exactly one solution in
\((0,1)\). At that probability every actor is indifferent and thus
every mixture is a best response to the opponents' choices; using
\(p^*\) by everyone is therefore an equilibrium. This proves existence
and uniqueness among symmetric independent mixed profiles, without claiming
uniqueness among all profiles. Expected utility at indifference is \(U-K\).

For an arbitrary independent profile, failure means all actors abstain;
its probability is \(\prod_i(1-p_i)\). Thus
\(P=1-\prod_i(1-p_i)\), and charging actor \(i\)'s cost with
probability \(p_i\) gives \(E_i=UP-Kp_i\). At the symmetric equilibrium,
\[
 P^*=1-(1-p^*)^N=1-(K/U)^{N/(N-1)}.
\]
Set \(r=K/U\in(0,1)\). For positive exponents, \(r^t\) strictly
decreases as \(t\) increases, since \(\log r<0\). As integer
\(N\geq2\) increases, both \(1/(N-1)\) and
\(N/(N-1)=1+1/(N-1)\) decrease. Therefore
\(r^{1/(N-1)}\) and \(r^{N/(N-1)}\) increase, and their complements
\(p^*\) and \(P^*\) decrease. The exponents tend to zero and one,
respectively, so the limits are zero and \(1-r\).

For contrast, if \(p\in(0,1)\) is held fixed rather than adjusted
to equilibrium, \(1-(1-p)^N\) strictly increases with \(N\).
This verifies the source's opposing-mechanism argument and identifies
which behavioral adjustment makes the equilibrium net effect negative.
\end{proof}
\end{annexitem}

\begin{annexitem}{annex:p17.rat104-cooperative-mixture}{thm:p17-rat104-cooperative-mixture}{Joint improvement, Pareto efficiency and correlation boundary}
\begin{proof}
At a symmetric independent profile with common probability \(p\), every
actor's expected utility is
\[
 F(p)=U\bigl[1-(1-p)^N\bigr]-Kp,
\qquad F'(p)=NU(1-p)^{N-1}-K.
\]
The derivative strictly decreases on \([0,1]\), is positive at zero
and negative at one. Its unique zero is
\[
                  s=1-(K/(NU))^{1/(N-1)}=p^{**}.
\]
Consequently \(F\) strictly increases up to \(s\) and strictly
decreases afterward; \(s\) is its unique maximizer. Since
\(0<K/(NU)<K/U<1\), it follows that \(0<p^*<s<1\).
Write \(a=1-s\), so \(U a^{N-1}=K/N\). Then
\[
\begin{aligned}
 F(s)&=U(1-a^N)-K(1-a)\\
     &=U-K+Ka-Ua^N\\
     &=U-K+\frac{N-1}{N}Ka=U-qK,\\
 q&=1-\frac{N-1}{N}a.
\end{aligned}
\]
Here \(0<a<1\) gives \(0<q<1\), and the displayed positive term
proves \(F(s)>U-K\).

The symmetric maximization alone would not prove efficiency over asymmetric
profiles. To establish that stronger independent-domain claim, take any
profile \((p_1,\ldots,p_N)\), put \(m=\max_i p_i\) and choose
\(j\) attaining the maximum. Since \(1-p_i\geq1-m\),
\[
 E_j=U\left[1-\prod_i(1-p_i)\right]-Km
       \leq U[1-(1-m)^N]-Km=F(m)\leq F(s).
\]
If every actor were at least as well off as at \((s,\ldots,s)\),
these inequalities would both be equalities. Uniqueness of the maximum
forces \(m=s\). Since \(1-s>0\), equality in the product inequality
forces \(p_i=s\) for every \(i\): increasing any one factor above
\(1-s\) would strictly increase the product. The entire profile is
therefore the original profile, and nobody can be strictly better off.
This proves Pareto efficiency over all independent profiles without
interpersonal utility summation.

Against \(N-1\) opponents each using \(s\), abstention yields
\[
        U[1-a^{N-1}]=U-K/N>U-K,
\]
the utility from contributing. An actor strictly prefers pure abstention,
so the symmetric profile at \(s>0\) is not an equilibrium.

For the correlated-domain boundary, choose exactly one actor uniformly
from the \(N\) actors and direct that actor alone to contribute.
Each actor receives the benefit with probability one and pays the cost
with probability \(1/N\), obtaining \(U-K/N\). Relative to the
independent symmetric profile this raises each utility by
\[
 (U-K/N)-F(s)=\frac{N-1}{N}K(1-a)
                          =\frac{N-1}{N}Ks>0.
\]
The resulting joint distribution is correlated, not an independent profile.
The calculation establishes a feasible payoff comparison in the enlarged
domain, not voluntary compliance with the proposed recommendations.
\end{proof}
\end{annexitem}

\begin{annexitem}{annex:p17.rat104-provision-correction}{thm:p17-rat104-provision-correction}{Correcting the source's unconditional cooperative group-size claim}
\begin{proof}
Independence and \(1-p^{**}=(r/N)^{1/(N-1)}\) give
\[
              P^{**}(N)=1-(r/N)^{N/(N-1)}.
\]
For \(r=1/2\), at \(N=2\) the failure probability is \(1/16\).
At \(N=3\) it is \((1/6)^{3/2}\). Both are positive and
\[
       \bigl[(1/6)^{3/2}\bigr]^2=1/216>1/256=(1/16)^2.
\]
Thus the failure probability increases and provision decreases between
these allowed integer sizes. This disproves the unconditional increase
in note~5 of the source; it does not alter the formula for \(p^{**}\).

To characterize rather than merely refute the assertion, extend the
formula to real \(N>1\). Put
\[
 f(N)=\frac{N}{N-1}\log(r/N),\qquad P^{**}(N)=1-\exp(f(N)).
\]
Ordinary differentiation gives
\[
 f'(N)=\frac{\log(N/r)-(N-1)}{(N-1)^2},\qquad
 \frac{dP^{**}}{dN}=-\exp(f(N))f'(N).
\]
The exponential and squared denominator are positive, proving the stated
derivative criterion, including equality. For \(N\geq2\), the function
\(h(N)=N-1-\log(N/r)\) strictly increases because
\(h'(N)=1-1/N>0\), and tends to infinity. Hence the provision curve
eventually increases, which also yields eventual increase for the integer
sequence. If \(r\geq2/e\), then
\(h(2)=1-\log(2/r)\geq0\), and \(h(N)>0\) for \(N>2\).
So the integer sequence is strictly increasing for all \(N\geq2\),
including the endpoint case \(r=2/e\).

Finally, since \(r/N\in(0,1)\) and \(N/(N-1)>1\),
\[
              0<(r/N)^{N/(N-1)}<r/N\longrightarrow0.
\]
Thus \(P^{**}(N)\) tends to one. An eventual limit and derivative
criterion cannot be replaced by unconditional monotonicity at every finite
group size.
\end{proof}
\end{annexitem}
